\chapter*{Conclusion générale}
\addcontentsline{toc}{chapter}{Conclusion générale}

\section*{Rappel du contexte et de la démarche}

Ce rapport a présenté la conception et l'implémentation d'AGT Infra, une
plateforme interne de gestion d'infrastructure et de supervision
développée au sein d'AG Technologies. Le constat de départ était simple
mais concret : un parc de serveurs physiques hébergeant des applications
critiques, piloté de façon manuelle et dispersée, sans point de contrôle
unique ni traçabilité homogène des actions effectuées.

L'état de l'art a montré qu'aucune solution existante ne réunissait
simultanément l'ensemble des contraintes propres au contexte d'AG
Technologies, ce qui a justifié un développement interne. L'analyse et
la conception ont ensuite posé les fondations du système : une
architecture qui pilote des briques externes déjà éprouvées sans jamais
les réimplémenter, et un modèle de contrôle d'accès dynamique où un
refus explicite prévaut toujours sur une autorisation plus large.

\section*{Bilan du projet}

L'implémentation a confirmé la viabilité de cette conception. Sur les
quatorze domaines fonctionnels définis, la majorité est aujourd'hui
réalisée et vérifiée en conditions réelles :

\begin{itemize}
  \item gestion des serveurs et des conteneurs Docker ;
  \item déploiement continu, avec retour arrière automatique en cas
        d'échec ;
  \item publication applicative sécurisée, sous HTTPS ;
  \item supervision autonome par agent dédié ;
  \item contrôle d'accès entièrement piloté en base de données.
\end{itemize}

Certaines limites restent assumées, en particulier l'absence de
sauvegardes automatiques, de montée en charge automatique et d'un
moteur de sécurité actif, qui demeurent les axes les plus critiques pour
une éventuelle suite du projet.

\section*{Compétences techniques développées}

Ce stage a permis de mettre en pratique un cycle de développement
rigoureux (audit, plan, implémentation, test) sur un projet de taille
réelle, ainsi que de concevoir et manipuler des mécanismes qui dépassent
le cadre d'un exercice académique classique :

\begin{itemize}
  \item conception d'un système de contrôle d'accès à la fois flexible
        et sûr (RBAC/IBAC dynamique) ;
  \item orchestration de serveurs et de conteneurs à distance, sans
        agent installé sur les cibles ;
  \item mise en place de flux temps réel par WebSocket ;
  \item sécurisation d'évènements entrants par signature (webhooks) ;
  \item conception d'une architecture volontairement dépourvue de tâche
        de fond côté serveur.
\end{itemize}

\section*{Bilan personnel}

Au delà des aspects techniques, ce stage a été l'occasion de développer
des compétences méthodologiques et humaines, en particulier :

\begin{itemize}
  \item mieux découper un travail complexe en étapes maîtrisables,
        plutôt que de l'aborder d'un seul bloc ;
  \item adopter une pensée orientée besoin, en partant systématiquement
        de ce qu'un module doit accomplir avant de choisir comment le
        construire ;
  \item savoir remettre en question un choix déjà fait, lorsque
        l'expérience montre qu'il ne tient plus face au besoin réel ;
  \item être plus rigoureux sur la documentation, condition nécessaire
        pour reprendre un travail complexe après une interruption ;
  \item recourir à des maquettes pour se donner une vue d'ensemble,
        dès qu'un sujet devient trop complexe pour être abordé
        directement dans le code.
\end{itemize}

Travailler avec une équipe, au delà du travail individuel sur le code,
a également demandé une capacité de remise en question, une
compréhension approfondie des choix effectués (chaque décision devant
pouvoir être justifiée), ainsi qu'une rigueur de ponctualité dans le
suivi des échanges et des livraisons.

\section*{Apport pour le projet professionnel}

Ce stage a confirmé et affiné mon orientation professionnelle. Avant
cette expérience, je me dirigeais vers un profil de développeur backend
et SecOps. Ce projet m'a permis de découvrir le \textit{platform
engineering}, une discipline qui prolonge naturellement le DevOps en
construisant des plateformes internes destinées à outiller d'autres
équipes, et que je considère désormais comme une direction d'avenir pour
ce domaine. Cette découverte oriente aujourd'hui plus précisément la
suite de mon parcours professionnel.